\chapter[Introdução]{Introdução}

\textcolor{black}{
A robótica móvel autônoma tem avançado de forma significativa nas últimas décadas,
impulsionada pela crescente demanda por sistemas capazes de operar em ambientes
desconhecidos sem intervenção humana.
Nesse contexto, os robôs exploradores de labirintos, popularmente conhecidos como
\textit{micromouse}, constituem uma plataforma clássica e amplamente estudada para
o desenvolvimento e validação de algoritmos de navegação autônoma, mapeamento e
tomada de decisão em tempo real \cite{387215}.
Em trabalho publicado no \textit{IEEE Transactions on Education}, foi reportada a
integração bem-sucedida do projeto \textit{micromouse} na grade curricular de cursos
de graduação da California State University, Fullerton, demonstrando que o
desenvolvimento de um robô autônomo capaz de encontrar um destino predeterminado
em um labirinto desconhecido constitui uma experiência pedagógica eficaz para
capacitar estudantes nas áreas de hardware, controle, sensoriamento, integração de
sistemas e algoritmos de planejamento de rotas.
}

A competição \textit{Micromouse}, regulamentada pelo
\textit{Institute of Electrical and Electronics Engineers} (IEEE),
é realizada anualmente em diversas regiões do mundo, incluindo América do Norte,
Europa e Ásia, e estabelece como desafio central a navegação autônoma de um robô
de pequeno porte por labirintos de dimensões padronizadas, sem qualquer intervenção
externa durante o percurso.
A norma IEEE vigente define, entre outros requisitos, que o robô não deve exceder
25~cm em comprimento ou largura, que a locomoção por meios aéreos ou por propulsão
a combustão é proibida, e que o código embarcado não pode ser alterado durante a
execução do trajeto~\footnote{
    IEEE Region 2 Student Activities Committee.
    \textit{MicroMouse Rules 2020}.
    Disponível em: \url{https://attend.ieee.org/r2sac-2020/wp-content/uploads/sites/175/2020/01/MicroMouse_Rules_2020.pdf}.
    Acesso em: maio 2026.
}.

Do ponto de vista científico, o problema de resolução de labirintos envolve subáreas
consolidadas da inteligência artificial e da engenharia de controle.
O algoritmo \textit{Flood Fill}, amplamente empregado em implementações de
\textit{micromouse}, permite que o robô construa um mapa incremental do labirinto
e calcule o caminho de menor custo até o objetivo com base nas informações sensoriais
coletadas em tempo real.
Conforme destacado por \cite{387215}, o sucesso no desenvolvimento de um
\textit{micromouse} competitivo exige não apenas domínio técnico individual,
mas também trabalho em equipe, gestão de tempo e capacidade de comunicação ---
habilidades centrais na formação de engenheiros e diretamente alinhadas às
competências desenvolvidas ao longo do presente projeto integrador.

No que diz respeito à legislação e normatização pertinente, o projeto está sujeito
às diretrizes do IEEE para competições de \textit{Micromouse}, às normas da
Associação Brasileira de Normas Técnicas (ABNT) aplicáveis a equipamentos
eletroeletrônicos de baixa tensão --- em especial a ABNT NBR IEC 60950-1 ---,
e às regulamentações da Agência Nacional de Telecomunicações (ANATEL) referentes
à utilização de módulos de comunicação sem fio em território nacional, conforme a
Resolução ANATEL n\textordmasculine{}~715, de 23 de setembro de 2019, que aprova
o Regulamento sobre Equipamentos de Radiocomunicação de Radiação Restrita.

O mercado global de robótica educacional e de competição tem apresentado crescimento
consistente, com estimativas apontando para um valor superior a US\$~1,8~bilhão em
2023 e projeção de expansão a taxas superiores a 12\% ao ano até 2030, segundo
relatório da \textit{Markets and Markets Research} (2023).
Apesar disso, soluções comerciais completas voltadas especificamente para
\textit{micromouse} permanecem escassas e de elevado custo de importação, criando
uma lacuna relevante para o desenvolvimento de plataformas abertas, acessíveis e
adaptadas à realidade brasileira.

Diante desse cenário, o presente projeto propõe o desenvolvimento de um sistema
\textit{micromouse} autônomo de baixo custo, concebido a partir de componentes
amplamente disponíveis no mercado nacional e projetado para atender integralmente
aos requisitos da competição IEEE.
A solução integra hardware embarcado, algoritmos de navegação e mapeamento, e um
sistema web de telemetria em tempo real, constituindo uma plataforma didática e
funcional que, assim como reportado por \cite{387215}, demonstra a viabilidade de
sistemas robóticos autônomos desenvolvidos no contexto acadêmico, contribuindo
tanto para a formação multidisciplinar de engenheiros quanto para a democratização
da robótica de competição no Brasil.
